\odiseachapter{voz}{La voz de Circe}\label{cap:voz}

Todo un año de vacaciones pagadas, cuando, supuestamente, los aqueos están deseando volver a su patria. ¿Se les ha quitado el ansia gracias a la magia de la hechicera, o es a su capitán, alojado en la cama relinda de Circe, al que se le han pasado las prisas? Más bien parece lo segundo, ya que es la tripulación la que interpela a Odiseo:

\begin{versocita}
«Hombre de dios, del solar de la patria tendrás que acordarte\\
ya si la voz de los cielos dictado dejó que te salves\\
y que a tu casa techalta al fin llegues y al lar de los padres».
\end{versocita}

Ulises, por fin, recuerda que lo suyo es volver a casa y dejarse de juergas.

\begin{versocita}
Yo subí entonces de nuevo a la cama relinda de Circe\\
y a sus rodillas le supliqué, y la diosa a mí oíame,\\
y, levantando de verbos alada la voz, por fin dije:
\stanzabreak
«Cúmpleme, Circe, ahora eso de lo que me hicieras promesa,\\
mandarme a casa; que el alma se me anda por eso ya inquieta,\\
y la de mis compañeros, que mi corazón atormentan\\
aquí y allá lamentándose en cuanto que tú no estás cerca».
\end{versocita}

La diosa parece consentir, pero pone una condición tremenda:

\begin{versocita}
«Sangre Laercia, Odiseo milmañas, hijo del cielo,\\
no os quedéis ya por más tiempo en mi casa si ya no hay deseo.\\
Mas fuerza es que otro viaje cumpláis y destino primero\\
adonde tienen la horrenda Perséfone y Hades su reino\\
por del aliento tener del tebano Tiresias agüero,\\
aún con el seso en su sitio, que solo a este ciego agorero,\\
cuando murió, le mantuvo Perséfone su entendimiento\\
y lucidez, en un mundo de sombras en revoloteo».
\end{versocita}

%Terrible sitio el Hades, del que hablaremos enseguida. 
%
%%A Ulises, claro, la propuesta no le hace ninguna gracia:
%
%\begin{versocita}
%Eso me dijo, y a mí el corazón se me desmoronaba.\\
%Y allí sentado en el lecho me puse a llorar y mi alma\\
%ya no quería vivir ni quería del sol ver más llama.\\
%Mas cuando de revolverme y llorar se colmaron las ganas,\\
%tuve palabras entonces para contestar su demanda:
%\stanzabreak
%«Circe, ¿y quién va a conducirnos ahora por ese camino?\\
%Nadie ha llegado a la casa de Hades en negro navío».
%\end{versocita}
%
%La hechicera le explica que será un viento mágico el que lo lleve.
%
%\begin{versocita}
%«Hijo del cielo, Odiseo milmañas, sangre Laercia,\\
%pierde cuidado, por falta de guía en la nave no hay pena;\\
%fíjale el palo sin más y que blancas le caigan las velas\\
%y espera solo: verás cómo el soplo del Bóreas la leva.\\
%Solo después de cruzar el Océano con tu goleta,\\
%en punta baja hallarás los bosques de Persefonea,\\
%álamos altos y sauces machíos entre cañaveras,\\
%junto al Océano remolinhondo allí el barco fondea y\\
%de la mansión aguanosa de Hades darás en la senda».
%\end{versocita}
%
%Le explica también lo que tiene que hacer al llegar:
%
%\begin{versocita}
%«Al Aqueronte el Piriflegetonte sus rápidos echa\\
%con el Cocito, que es brazo de Estigia la malagunera,\\
%y hay una peña donde ambos dos ríos soneros se encuentran.
%\stanzabreak
%A ella tú acércate, héroe, y ---escúchame y ripio no pierdas---\\
%abre de un codo a lo largo y de un codo a lo hondo una cuenca\\
%para que viertas por todos los muertos en ella una ofrenda:\\
%primero de lechimel, dulce vino al seguido, y, sobre ellas,\\
%una tercera de agua; y harina blanca le avientas.
%\stanzabreak
%Invocarás de los muertos, de hinojos, las nulas cabezas,\\
%y les dirás que ya en Ítaca harás de una vaca horra y buena\\
%un sacrificio en tu casa y el lar llenarás de preseas,\\
%y que un carnero, además, para él solo y aparte, a Tiresias\\
%le inmolarás, todo negro, de vuestro rebaño excelencia.
%\stanzabreak
%Cuando a las ínclitas razas de muertos rezadas las tengas,\\
%degollarás un cordero tú allí y una oveja bien negra\\
%con la testuz inclinada al Erebo; tú vuelve tu testa\\
%como a seguir las corrientes del río: allí ya la cuenta\\
%vas a perder de las sombras de muertos difuntos que llegan.
%\stanzabreak
%En ese mismo momento a tus hombres apremia y ordena\\
%que, desollando las reses que a crudo bronce en la tierra\\
%tú degollaste, las quemen, rezándole al dios y a la dea,\\
%a Hades tremendo y con él a la hórrida Persefonea.
%\stanzabreak
%Y tú, de la hoja tirando que llevas al muslo sujeta,\\
%mira que muerto ninguno acerque su huera cabeza\\
%a aquella sangre hasta que de Tiresias te venga respuesta.\\
%Presto, caudillo de hombres, allí acudirá ese profeta\\
%que te dirá el buen camino y el tiempo de tu carrera\\
%y ese regreso que por la llanura de peces tú tengas».
%\end{versocita}
%
Odiseo se resigna a hacer el viaje y avisa a sus hombres. Pero no van a marcharse de la isla de Circe sin sufrir otra desgracia:

\begin{versocita}
Eso les dije, y su aliento y arrojo a mi voz se rendían.\\
Mas ni de allí pude ilesos sacar a mis hombres. Había\\
uno, el más joven de todos, Elpénor, que ni era de hombría\\
para la guerra excesivo ni por su cordura lucía,\\
que por el santo palacio de Circe y sin compañía,\\
fresca buscando y pesado de vino, a dormir se fue arriba.\\
Cuando las voces oyó y el trajín, porque ya se movían\\
sus compañeros, saltó de repente sin ver que debía\\
ir hacia atrás y bajar por la larga escalera: de crisma\\
desde el tejado cayó y su cuello quedó por la espina\\
roto y quebrado, y al Hades se hundió su resuello de vida.
\end{versocita}

Una muerte absurda, que conmueve por lo arbitrario, lo casual, lo fuera de lugar. Parece que Homero quiera recordarnos que la parca nos espera detrás de cada esquina. El pobre Elpénor, el más joven de los compañeros, aparece en escena solo para que las furias lo arrojen, de un zarpazo, al Hades.

El canto termina de manera misteriosa, con Circe paseándose invisible por la nave de los aqueos:

\begin{versocita}
Cuando a la léveda nave, a la orilla del mar, afligidos,\\
lágrima en flor derramando hacíamos todos camino,\\
Circe nos vino a pasar por delante y al negro navío\\
ató una oveja bien negra con un carnero chiquito, y\\
volvió a pasar sin ser vista de largo: a dios que no quiso\\
¿quién con sus ojos lo vio alguna vez cuando fue y cuando vino?
\end{versocita}

¿Quién es Circe y qué pasó exactamente en Eea? No hace falta que repita que Ulises no es un narrador fiable y es bien posible que la historia que les está contando a los feacios incluya omisiones, fantasías o fabulaciones. ¿Por qué la hechicera, después de intentar convertirlo en cerdo como a todos sus hombres, accede tan rápidamente a sus deseos? ¿De verdad está asustada por las amenazas de un mortal? ¿Tiene miedo cuando le abraza las rodillas, o está jugando con él por divertirse? Ciertamente Odiseo no se hace de rogar para subir con ella al lecho, ni le parece tiempo excesivo pasarse un año en la isla. De hecho, son sus hombres los que tienen que suplicarle que sigan viaje.

Quizás hay otra lectura de Circe. En el epílogo de \emph{Abandonando Ítaca}, Candelas Gala escribe:

\begin{quotation}
La elección de seguir la dirección de Circe rompe con las expectativas. Ulises renuncia a la seguridad del hogar en Ítaca, a la lealtad y fidelidad de una esposa, por una situación tan potencialmente sorprendente e inesperada como el cambio del mar. Penélope le servirá un tejido/texto de hilos predecibles cuya linealidad acabará en la muerte; pero los nudos y giros de Circe le llevarán por caminos tortuosos, impredecibles, incluso peligrosos, pero vivos, donde lo inesperado contiene una plenitud vital que implica azar y probabilidad, pero también creatividad permanente.
\end{quotation}

En mi versión del poema, Ulises regresa a Ítaca, como se espera de él, pero es incapaz de olvidar a la hechicera. Muchos años más tarde, desde su todavía ardiente senectud, escribe:

\begin{versopropio}
Ha pasado mucho tiempo.\\
Intentas invocar su rostro, y lo único que puedes recordar\\
son sus ojos ardientes y su ávida boca.\\
Nunca ha hablado mucho, esa bruja,\\
ni una sola vez dijo\\
que te amaba.
\end{versopropio}

En realidad, tampoco se lo dijo nunca Penélope. Una cosa es el sexo y otra declarar los sentimientos. Nadie en la \emph{Odisea} habla nunca de amor.

Pero el recuerdo más lacerante para el anciano es la voz de Circe:

\begin{versopropio}
A menudo, la recuerdas cantando,\\
como aquella vez, cuando cantó en el entierro de uno de tus hombres,\\
era un don nadie, un pobre soldado que una noche bebió demasiado\\
y se subió al tejado, quizás en un vano intento\\
de ver desde más cerca las estrellas, perdió el pie\\
y se mató al caer, lo encontrasteis más tarde,\\
vagando por el Hades, la sombra de una sombra,\\
no más insignificante en la muerte de lo que fue en su vida.
\stanzabreak
Y sin embargo, él también, una vez, estuvo vivo.
\stanzabreak
Circe lo vio, ella que no tenía paciencia con los mortales,\\
y aún menos con los hombres,\\
se sintió, sin embargo, conmovida por el hecho ineludible\\
de que el pobre muchacho, tendido en su mesa, con el cuello roto\\
y el asombro en sus ojos, como si en su último momento\\
hubiese captado una verdadera visión de las constelaciones,\\
aquel niño, que yacía para siempre inmóvil, aquella cosa cuya forma\\
no seguiría siendo humana mucho tiempo, ayer mismo,\\
estuvo respirando bajo el Sol.
\stanzabreak
Y cuando empezó a cantar, su voz era tan pura,\\
y se alzó hasta las esferas,\\
cantando en la noche, como si nadie antes de ella\\
hubiese cantado, o llorado.
\end{versopropio}

Circe, en el recuerdo del viejo Ulises, deja de ser una diosa altiva, distante y no poco malvada, para convertirse en una mujer capaz de llorar por una muerte absurda:

\begin{versopropio}
Lloró por el joven que yacía muerto a sus pies,\\
y por todos los jóvenes que perecieron en Troya,\\
o que fueron borrados por Cíclopes y Lestrigones,\\
por sirenas y tormentas, cortados, quemados, hundidos,\\
con sus miembros perfectos cruelmente mutilados,\\
y la luz, para siempre robada de sus ojos.
\stanzabreak
Lloró también por los ancianos, como tu padre,\\
cuya vida ha sido larga,
\stanzabreak
aunque la vida humana nunca es bastante larga.\\
Los viejos, que van siendo niños de nuevo\\
pues todos los humanos, al envejecer,\\
regresan a su infancia.
\stanzabreak
Y como ocurre a un niño, el viejo está tan asustado\\
como el muchacho que masacraste aquella noche en Troya,\\
tu espada atravesó su vientre, sus tripas se desparramaron,\\
y trató de sostenerlas con sus manos, y llamó a su madre,\\
no sentiste piedad por él entonces, pero ahora está aquí,\\
contigo (y otros tantos fantasmas),\\
mientras te acuerdas de Circe.
\end{versopropio}

Ulises, consumido por remordimientos tardíos, se aferra al recuerdo de aquella voz prístina como se aferraba al madero que lo llevó a la isla de Nausícaa.

\begin{versopropio}
Sí, ella cantó. Sus ojos estaban cerrados, su pecho en alto.\\
En su voz podías escuchar las cuerdas de la lira\\
que tocaba Orfeo. Y la canción también lamentaba a Eurídice,\\
que no pudo ser rescatada del inframundo,\\
y se lamentó por la esposa de Héctor, esclavizada,\\
y se lamentó por todas esas otras mujeres que fueron tomadas,\\
por carniceros como tú.
\end{versopropio}

Y es esa voz la que le trae claridad, la que le desvela quién es, quién ha sido siempre:

\begin{versopropio}
La escuchaste cantar, y su voz era hermosa y, a un tiempo,\\
un dedo acusador, exigiendo saber por qué.\\
La escuchaste cantar, preguntándote\\
cómo es que un monstruo como tú merecía tanta hermosura.\\
Un monstruo como tú, que no debería tener corazón,\\
sino una fría piedra bombeando su negra sangre, y sin embargo,\\
abrumado por la pena, apenas podías respirar.
\stanzabreak
Sí, ella cantó. Su voz era un espejo, donde veías tus crímenes,\\
pero también una oración compasiva,\\
y un ofrecimiento de perdón. Su voz dijo\\
monstruo o no, eres solo un hombre, y todos los hombres,
\stanzabreak
lo merezcan o no,\\
pueden ser redimidos.
\end{versopropio}
